\section{Case Study 1: DeepSeek R1}

DeepSeek R1 是课堂里第一个重点案例。它之所以重要，不只是因为性能强，还因为它把一条相对简洁的 RLVR recipe 推到了大众视野：\textbf{强 verification reward + GRPO + 合理的 SFT 初始化 + 后续蒸馏}。

\begin{figure}[H]
\centering
\includegraphics[width=0.95\textwidth]{slides-images/slide-037.jpg}
\caption{案例总览：DeepSeek R1、Kimi K1.5、Qwen 3}
\end{figure}

\begin{figure}[H]
\centering
\includegraphics[width=0.95\textwidth]{slides-images/slide-038.jpg}
\caption{DeepSeek R1 的重要性：性能强、公开 recipe 相对简单、改变了外界对 reasoning training 的预期}
\end{figure}

\subsection{R1-zero 和 R1}

课堂区分了两个版本：
\begin{itemize}
    \item \textbf{R1-zero}：更“纯”的 RLVR 设置，用 accuracy reward 和 format reward 推动训练。
    \item \textbf{R1}：在 R1-zero 的基础上，再加 SFT 初始化、语言一致性 reward、以及更实用的后处理阶段。
\end{itemize}

R1-zero 的有趣现象是：训练中会出现更长的 CoT，有时还会冒出所谓的 `aha moment`。但课上也提醒，不能把这些现象讲得太神秘，因为它们也可能部分来自目标函数偏置和基座模型本身的能力。

\begin{figure}[H]
\centering
\includegraphics[width=0.95\textwidth]{slides-images/slide-040.jpg}
\caption{R1-zero 的受控实验：accuracy reward + format reward}
\end{figure}

\begin{figure}[H]
\centering
\includegraphics[width=0.95\textwidth]{slides-images/slide-041.jpg}
\caption{R1-zero 训练里的现象：更长的 CoT 和所谓的 `aha moment`}
\end{figure}

\begin{figure}[H]
\centering
\includegraphics[width=0.95\textwidth]{slides-images/slide-043.jpg}
\caption{R1 的关键差异：SFT 初始化、语言一致性 reward、以及后续的非可验证任务阶段}
\end{figure}

\subsection{蒸馏和失败路线}

DeepSeek R1 还有一个很重要的结论：\textbf{reasoning 能蒸馏}。也就是说，让大模型先用 RLVR 形成 reasoning traces，再把这些长 CoT 蒸馏给较小模型，是一条非常实用的路线。

同时，课堂也提醒了另一个事实：很多大家曾经寄予厚望的路线，未必是最后赢家，例如 PRM、MCTS 等。这个领域里真正起作用的，往往不是“理论上最花哨的东西”，而是能稳定跑通的 recipe。

\begin{figure}[H]
\centering
\includegraphics[width=0.95\textwidth]{slides-images/slide-049.jpg}
\caption{蒸馏：用 R1 生成 CoT traces，再教给较小的 student model}
\end{figure}

\begin{figure}[H]
\centering
\includegraphics[width=0.95\textwidth]{slides-images/slide-050.jpg}
\caption{一些没有最终跑通的路线：PRM、MCTS 等}
\end{figure}

